\textbf{Phần 1. Trắc nghiệm lựa chọn}
\setcounterpageref{ex}{0}

\begin{ex} %[5D1N2]
Kết quả của phép tính $7\,318 + 3\,191$ là
\choice
{$10\,409$}
{$10\,519$}
{$10\,499$}
{\True $10\,509$}
\loigiai{
Đặt tính và tính: $7\,318 + 3\,191 = 10\,509$.\\
\textbf{Đáp án của câu hỏi là D}
}
\end{ex}

\begin{ex} %[5D1N2]
Kết quả của phép trừ $83\,500 - 28\,150$ là
\choice
{\True $55\,350$}
{$55\,450$}
{$65\,350$}
{$55\,250$}
\loigiai{
Thực hiện phép trừ: $83\,500 - 28\,150 = 55\,350$.\\
\textbf{Đáp án của câu hỏi là A}
}
\end{ex}

\begin{ex} %[5D1H2]
Kết quả của phép nhân $307 \times 15$ là
\choice
{$4\,505$}
{$4\,595$}
{\True $4\,605$}
{$4\,615$}
\loigiai{
Thực hiện phép nhân: $307 \times 15 = 4\,605$.\\
\textbf{Đáp án của câu hỏi là C}
}
\end{ex}

\begin{ex} %[5D1H2]
Thương của phép chia $4\,488 : 34$ là
\choice
{$122$}
{$142$}
{$134$}
{\True $132$}
\loigiai{
Đặt tính chia: $4\,488 : 34 = 132$.\\
\textbf{Đáp án của câu hỏi là D}
}
\end{ex}

\begin{ex} %[5D1H2]
Biểu thức nào dưới đây có giá trị bằng biểu thức $2 \times 75 + 2 \times 25$?
\choice
{\True $2 \times (75 + 25)$}
{$2 \times (75 - 25)$}
{$2 + (75 \times 25)$}
{$(2 + 2) \times (75 + 25)$}
\loigiai{
Áp dụng tính chất nhân một số với một tổng: $2 \times 75 + 2 \times 25 = 2 \times (75 + 25)$.\\
\textbf{Đáp án của câu hỏi là A}
}
\end{ex}

\begin{ex} %[5D1H2]
Giá trị của biểu thức $3\,713 - 200 \times 5$ là
\choice
{$17\,565$}
{$2\,613$}
{\True $2\,713$}
{$3\,513$}
\loigiai{
Thực hiện phép nhân trước: $200 \times 5 = 1\,000$. Sau đó trừ: $3\,713 - 1\,000 = 2\,713$.\\
\textbf{Đáp án của câu hỏi là C}
}
\end{ex}

\begin{ex} %[5D1H2]
Giá trị của biểu thức $1\,500 + (750 + 250) : 2$ là
\choice
{$1\,000$}
{\True $2\,000$}
{$2\,500$}
{$1\,750$}
\loigiai{
Thực hiện trong ngoặc trước: $750 + 250 = 1\,000$.\\
Sau đó chia: $1\,000 : 2 = 500$.\\
Cuối cùng cộng: $1\,500 + 500 = 2\,000$.\\
\textbf{Đáp án của câu hỏi là B}
}
\end{ex}

\begin{ex} %[5D1H2]
Tính nhanh biểu thức $25 \times 99 \times 4$ ta được kết quả là
\choice
{\True $9\,900$}
{$990$}
{$99\,000$}
{$990\,000$}
\loigiai{
Sử dụng tính chất giao hoán và kết hợp: $(25 \times 4) \times 99 = 100 \times 99 = 9\,900$.\\
\textbf{Đáp án của câu hỏi là A}
}
\end{ex}

\begin{ex} %[5D1H2]
Tính giá trị của biểu thức $2\,025 \times 17 + 83 \times 2\,025$ bằng cách thuận tiện nhất:
\choice
{$202\,500\,000$}
{$2\,025\,000$}
{\True $202\,500$}
{$20\,250$}
\loigiai{
Áp dụng tính chất phân phối: $2\,025 \times (17 + 83) = 2\,025 \times 100 = 202\,500$.\\
\textbf{Đáp án của câu hỏi là C}
}
\end{ex}

\begin{ex} %[5D1V2]
Mai mua hai gói bim bim hết số tiền là $18\,000$ đồng. Biết gói bim bim cua đắt hơn gói bim bim mực là $4\,000$ đồng. Giá tiền của gói bim bim mực là
\choice
{$11\,000$ đồng}
{\True $7\,000$ đồng}
{$14\,000$ đồng}
{$9\,000$ đồng}
\loigiai{
Đây là bài toán tìm hai số khi biết tổng và hiệu:\\
Giá tiền gói bim bim mực (số bé) là: $(18\,000 - 4\,000) : 2 = 7\,000$ đồng.\\
\textbf{Đáp án của câu hỏi là B}
}
\end{ex}

\begin{ex} %[5D1V2]
Rô-bốt bán bốn bức tranh với giá tiền lần lượt là: $85\,500$ đồng, $150\,000$ đồng, $425\,000$ đồng và $55\,500$ đồng. Trung bình mỗi bức tranh Rô-bốt bán được giá bao nhiêu tiền?
\choice
{$178\,000$ đồng}
{$180\,000$ đồng}
{$189\,000$ đồng}
{\True $179\,000$ đồng}
\loigiai{
Tổng số tiền bán $4$ bức tranh là:
\[
85\,500 + 150\,000 + 425\,000 + 55\,500 = 716\,000\text{ (đồng)}.
\]
Trung bình mỗi bức tranh có giá: $716\,000 : 4 = 179\,000$ đồng.\\
\textbf{Đáp án của câu hỏi là D}
}
\end{ex}

\begin{ex} %[5D1V2]
Bác Ba có $525\text{ kg}$ gạo chia đều vào các túi, mỗi túi $15\text{ kg}$. Hỏi $7$ túi gạo như vậy chứa tất cả bao nhiêu ki-lô-gam gạo?
\choice
{$75\text{ kg}$}
{\True $105\text{ kg}$}
{$120\text{ kg}$}
{$150\text{ kg}$}
\loigiai{
Số ki-lô-gam gạo trong $7$ túi là: $15 \times 7 = 105\text{ (kg)}$.\\
\textbf{Đáp án của câu hỏi là B}
}
\end{ex}

\textbf{Phần 2. Trắc nghiệm đúng sai}
\setcounterpageref{ex}{0}

\begin{ex} %[5D1H2]
Cho các biểu thức: $A = 25 \times 99 \times 4$, $B = 2\,025 \times 17 + 83 \times 2\,025$ và $C = 3\,713 - 200 \times 5$. Xét tính đúng sai của các khẳng định sau:
\choiceTF
{\True Giá trị của biểu thức $A$ bằng $9\,900$.}
{Giá trị của biểu thức $B$ là $20\,250$.}
{\True Biểu thức $C$ có giá trị là $2\,713$.}
{\True Tổng giá trị của biểu thức $A$ và biểu thức $C$ là $12\,613$.}
\loigiai{
\begin{itemize}
    \item $A = (25 \times 4) \times 99 = 100 \times 99 = 9\,900$.
    \item $B = 2\,025 \times (17 + 83) = 2\,025 \times 100 = 202\,500 \neq 20\,250$.
    \item $C = 3\,713 - 1\,000 = 2\,713$.
    \item $A + C = 9\,900 + 2\,713 = 12\,613$.
\end{itemize}
\begin{itemize}
    \item a) \textbf{Đúng}.
    \item b) \textbf{Sai}: $B = 202\,500$.
    \item c) \textbf{Đúng}.
    \item d) \textbf{Đúng}.
\end{itemize}
}
\end{ex}

\begin{ex} %[5D1V2]
Bác Ba chia đều $525\text{ kg}$ gạo vào các túi nhỏ, mỗi túi chứa $15\text{ kg}$. Bác bán mỗi túi gạo đó với giá $250\,000$ đồng. Xét tính đúng sai của các khẳng định sau:
\choiceTF
{\True Bác Ba chia được tất cả $35$ túi gạo.}
{$7$ túi gạo như vậy chứa $115\text{ kg}$ gạo.}
{\True Tổng số tiền bác Ba thu được khi bán hết số gạo là $8\,750\,000$ đồng.}
{Nếu bác Ba chỉ bán $20$ túi gạo thì số tiền thu được là $4\,500\,000$ đồng.}
\loigiai{
\begin{itemize}
    \item Tổng số túi gạo bác Ba chia được là: $525 : 15 = 35$ (túi).
    \item $7$ túi gạo chứa số ki-lô-gam là: $15 \times 7 = 105\text{ kg} \neq 115\text{ kg}$.
    \item Bán hết $35$ túi thu được: $35 \times 250\,000 = 8\,750\,000$ (đồng).
    \item Bán $20$ túi thu được: $20 \times 250\,000 = 5\,000\,000\text{ (đồng)} \neq 4\,500\,000$ đồng.
\end{itemize}
\begin{itemize}
    \item a) \textbf{Đúng}.
    \item b) \textbf{Sai}: $7$ túi chứa $105\text{ kg}$.
    \item c) \textbf{Đúng}.
    \item d) \textbf{Sai}: Thu được $5\,000\,000$ đồng.
\end{itemize}
}
\end{ex}

\textbf{Phần 3. Trả lời ngắn}
\setcounterpageref{ex}{0}

\begin{ex} %[5D1H2]
Tính tổng tất cả các số sau bằng cách thuận tiện: $400, 500, 300, 600, 280, 510, 100, 500, 490, 900, 720, 700$.
\shortans{6000}
\loigiai{
\textbf{Đáp số:} $6000$\\
Nhóm các cặp số tròn nghìn và tròn trăm:
\begin{align*}
S &= (400 + 600) + (500 + 500) + (300 + 700) + (100 + 900) + (280 + 720) + (510 + 490) \\
&= 1\,000 + 1\,000 + 1\,000 + 1\,000 + 1\,000 + 1\,000 = 6\,000.
\end{align*}
}
\end{ex}

\begin{ex} %[5D1H2]
Tìm chữ số thích hợp thay cho dấu hỏi chấm trong phép trừ sau:
\[
\begin{array}{r@{}l@{\quad}l}
  & 6\,0\,6\,1\,8\,? \\
- & 4\,3\,?\,1\,4\,9 \\
\hline
  & 1\,7\,1\,?\,3\,3
\end{array}
\]
Biết các dấu hỏi chấm từ trên xuống dưới, từ trái sang phải lần lượt là các chữ số. Hãy cho biết chữ số ở hàng đơn vị của số bị trừ (dấu hỏi chấm ở dòng thứ nhất).
\shortans{2}
\loigiai{
\textbf{Đáp số:} $2$\\
Ở hàng đơn vị: $? - 9 = 3 \Rightarrow ? = 12$ (viết $3$ nhớ $1$), vậy chữ số cần tìm là $2$.
}
\end{ex}

\begin{ex} %[5D1V2]
Mai mua hai gói bim bim hết số tiền là $18\,000$ đồng. Trong đó gói bim bim cua đắt hơn gói bim bim mực là $4\,000$ đồng. Hỏi giá tiền gói bim bim cua là bao nhiêu đồng?
\shortans{11000}
\loigiai{
\textbf{Đáp số:} $11000$\\
Giá tiền gói bim bim cua là: $(18\,000 + 4\,000) : 2 = 11\,000$ (đồng).
}
\end{ex}

\begin{ex} %[5D1H2]
Bác Ba chia đều $525\text{ kg}$ gạo vào các túi, mỗi túi $15\text{ kg}$. Hỏi bác Ba chia được tất cả bao nhiêu túi gạo?
\shortans{35}
\loigiai{
\textbf{Đáp số:} $35$\\
Số túi gạo chia được là: $525 : 15 = 35$ (túi).
}
\end{ex}

\begin{ex} %[5D1V2]
Bác Ba chia đều $525\text{ kg}$ gạo vào các túi, mỗi túi $15\text{ kg}$ và bán mỗi túi với giá $250\,000$ đồng. Hỏi sau khi bán hết số gạo đó, bác Ba thu được tất cả bao nhiêu nghìn đồng?
\shortans{8750}
\loigiai{
\textbf{Đáp số:} $8750$\\
Số túi gạo là $525 : 15 = 35$ túi. Số tiền thu được là $35 \times 250\,000 = 8\,750\,000$ đồng, tức $8\,750$ nghìn đồng.
}
\end{ex}

\begin{ex} %[5D1V2]
Rô-bốt bán bốn bức tranh với các mức giá: $85\,500$ đồng, $150\,000$ đồng, $425\,000$ đồng và $55\,500$ đồng. Hỏi trung bình mỗi bức tranh có giá bao nhiêu nghìn đồng?
\shortans{179}
\loigiai{
\textbf{Đáp số:} $179$\\
Tổng số tiền bán $4$ bức tranh là: $85\,500 + 150\,000 + 425\,000 + 55\,500 = 716\,000$ đồng.\\
Trung bình mỗi bức tranh có giá: $716\,000 : 4 = 179\,000$ đồng, tức $179$ nghìn đồng.
}
\end{ex}
